
\documentclass[aspectratio=169, 11pt]{beamer}

% ─── PAQUETES ─────────────────────────────────────────────────────────────────
\usepackage[T1]{fontenc}
\usepackage[utf8]{inputenc}
\usepackage[spanish]{babel}
\usepackage{lmodern}
\usepackage{booktabs}
\usepackage{graphicx}
\usepackage{xcolor}
\usepackage{tikz}
\usepackage{enumitem}
\usepackage{array}
\usepackage{colortbl}
\usepackage{fontawesome5}

% ─── TEMA BASE ────────────────────────────────────────────────────────────────
\usetheme{metropolis}
\setbeamertemplate{navigation symbols}{}

% ─── PALETA DE COLORES ────────────────────────────────────────────────────────
\definecolor{PriRed}{RGB}{178, 32, 40}
\definecolor{DarkBg}{RGB}{18, 18, 22}
\definecolor{LightBg}{RGB}{245, 246, 250}
\definecolor{MidGray}{RGB}{90, 90, 100}
\definecolor{AccBlue}{RGB}{30, 111, 165}
\definecolor{AccGreen}{RGB}{46, 125, 79}
\definecolor{AccAmber}{RGB}{184, 134, 11}
\definecolor{AccPurple}{RGB}{107, 63, 166}
\definecolor{AccTeal}{RGB}{26, 122, 150}

% ─── COLORES BEAMER ──────────────────────────────────────────────────────────
\setbeamercolor{normal text}{fg=DarkBg, bg=LightBg}
\setbeamercolor{frametitle}{fg=white, bg=PriRed}
\setbeamercolor{title separator}{fg=PriRed}
\setbeamercolor{progress bar}{fg=PriRed, bg=PriRed!30}
\setbeamercolor{block title}{fg=white, bg=PriRed}
\setbeamercolor{block body}{fg=DarkBg, bg=PriRed!8}
\setbeamercolor{alerted text}{fg=PriRed}
\setbeamercolor{itemize item}{fg=PriRed}
\setbeamercolor{itemize subitem}{fg=MidGray}
\setbeamercolor{enumerate item}{fg=PriRed}
\setbeamertemplate{itemize item}{\small\textbullet}
\setbeamertemplate{itemize subitem}{\tiny\textbullet}

% ─── TIPOGRAFÍA ──────────────────────────────────────────────────────────────
\setbeamerfont{frametitle}{size=\large, series=\bfseries}
\setbeamerfont{title}{size=\LARGE, series=\bfseries}
\setbeamerfont{subtitle}{size=\normalsize}

% ─── METADATOS ───────────────────────────────────────────────────────────────
\title{Detección y Eliminación Inteligente de Sellos\\para Mejora del OCR en Documentos Oficiales}
\subtitle{Objetivos, Estado del Arte y Metodología}
\author{Manuel Cruz Garrote}
\institute{Universidad del Rosario · Visión por Computador 2026-I}
\date{Marzo 2026}

% ─── COMANDOS AUXILIARES ─────────────────────────────────────────────────────
\newcommand{\lib}[1]{\texttt{\small\color{AccBlue}#1}}
\newcommand{\rojo}[1]{\textcolor{PriRed}{#1}}
\newcommand{\azul}[1]{\textcolor{AccBlue}{#1}}
\newcommand{\verde}[1]{\textcolor{AccGreen}{#1}}
\newcommand{\badge}[2]{%
  \tikz[baseline=(n.base)]\node[fill=#1!20, draw=#1!50, rounded corners=3pt,
  inner xsep=4pt, inner ysep=2pt, font=\ttfamily\scriptsize, text=#1!80!black] (n) {#2};%
}

% ─── BLOQUE ROJO (Ruta A) ────────────────────────────────────────────────────
\newenvironment{redblock}[1]{%
  \setbeamercolor{block title}{fg=white, bg=PriRed}%
  \setbeamercolor{block body}{fg=DarkBg, bg=PriRed!10}%
  \begin{block}{#1}}{%
  \end{block}}

% ─── BLOQUE AZUL (Ruta B) ────────────────────────────────────────────────────
\newenvironment{blueblock}[1]{%
  \setbeamercolor{block title}{fg=white, bg=AccBlue}%
  \setbeamercolor{block body}{fg=DarkBg, bg=AccBlue!10}%
  \begin{block}{#1}}{%
  \end{block}}

% ─── BLOQUE VERDE ────────────────────────────────────────────────────────────
\newenvironment{greenblock}[1]{%
  \setbeamercolor{block title}{fg=white, bg=AccGreen}%
  \setbeamercolor{block body}{fg=DarkBg, bg=AccGreen!10}%
  \begin{block}{#1}}{%
  \end{block}}

% ══════════════════════════════════════════════════════════════════════════════
\begin{document}
% ══════════════════════════════════════════════════════════════════════════════

% ─── PORTADA ──────────────────────────────────────────────────────────────────
\begin{frame}[plain]
  \maketitle
\end{frame}

% ─── ÍNDICE ───────────────────────────────────────────────────────────────────
\begin{frame}{Contenido}
  \setbeamertemplate{section in toc}[sections numbered]
  \tableofcontents[hideallsubsections]
\end{frame}

% ══════════════════════════════════════════════════════════════════════════════
\section{Contexto y Problema}
% ══════════════════════════════════════════════════════════════════════════════

\begin{frame}{El Problema}
  \begin{columns}[T]
    \begin{column}{0.52\textwidth}
      Los sellos oficiales superpuestos en documentos escaneados \alert{degradan significativamente} el rendimiento de los motores OCR, generando errores en la región cubierta.
      \vspace{12pt}
      \begin{block}{Pregunta de investigación}
        ¿En qué medida la \textbf{eliminación inteligente de sellos} mediante inpainting generativo mejora la precisión de la extracción automática de texto (OCR) en documentos oficiales escaneados?
      \end{block}
    \end{column}
    \begin{column}{0.44\textwidth}
      \begin{alertblock}{Causa principal}
        Los píxeles del sello se mezclan con el texto subyacente, confundiendo los detectores de caracteres.
      \end{alertblock}
      \vspace{8pt}
      \begin{exampleblock}{Impacto medible}
        La región cubierta genera tasas de error \textbf{hasta 3× superiores} al resto del documento.
      \end{exampleblock}
    \end{column}
  \end{columns}
\end{frame}

% ══════════════════════════════════════════════════════════════════════════════
\section{Objetivos}
% ══════════════════════════════════════════════════════════════════════════════

\begin{frame}{Objetivo General}
  \begin{block}{Definición}
    Desarrollar un \textbf{pipeline de visión por computador} capaz de detectar sellos en documentos oficiales escaneados en español, eliminarlos preservando el texto subyacente, y demostrar cuantitativamente que esta limpieza \textbf{mejora la precisión del OCR}.
  \end{block}
  \vspace{14pt}
  \begin{columns}[c]
    \begin{column}{0.24\textwidth}
      \centering{\Large\bfseries\rojo{> 5\%}}\\[4pt]{\scriptsize Mejora mínima en CER}
    \end{column}
    \begin{column}{0.24\textwidth}
      \centering{\Large\bfseries\rojo{6}}\\[4pt]{\scriptsize Fases del pipeline}
    \end{column}
    \begin{column}{0.24\textwidth}
      \centering{\Large\bfseries\rojo{100}}\\[4pt]{\scriptsize Documentos benchmark SROIE}
    \end{column}
    \begin{column}{0.24\textwidth}
      \centering{\Large\bfseries\rojo{25}}\\[4pt]{\scriptsize Documentos revisión humana}
    \end{column}
  \end{columns}
\end{frame}

\begin{frame}{Objetivos Específicos}
  \begin{enumerate}[label=\textbf{\rojo{\arabic*.}}, leftmargin=1.8em, itemsep=10pt]
    \item \textbf{Construir un dataset mixto} (sintético y real) de documentos en español con sellos, incluyendo \textit{ground truth} controlado para evaluación reproducible.
    \item \textbf{Comparar dos enfoques de detección:} segmentación clásica \textbf{HSV} versus detección por deep learning con \textbf{YOLOv8}, evaluando con métricas de IoU.
    \item \textbf{Desarrollar un sistema híbrido de eliminación} que combine separación cromática (sellos de color) con inpainting generativo \textbf{LaMa} (sellos negros), preservando el texto subyacente.
    \item \textbf{Evaluar la mejora en OCR} comparando \textbf{CER} y \textbf{WER} antes y después del pipeline — automáticamente y con validación humana.
  \end{enumerate}
\end{frame}

% ══════════════════════════════════════════════════════════════════════════════
\section{Estado del Arte}
% ══════════════════════════════════════════════════════════════════════════════

\begin{frame}{Métricas de Evaluación}
  \begin{columns}[T]
    \begin{column}{0.48\textwidth}
      \begin{block}{CER — Character Error Rate}
        Tasa de error a nivel de \textbf{carácter} basada en distancia de Levenshtein. Estándar para OCR en texto denso.
        $$\text{CER} = \frac{S + D + I}{N} \times 100\%$$
      \end{block}
    \end{column}
    \begin{column}{0.48\textwidth}
      \begin{block}{WER — Word Error Rate}
        Tasa de error a nivel de \textbf{palabra}. Más sensible a errores que cambian el significado completo de un token.
        $$\text{WER} = \frac{S + D + I}{N} \times 100\%$$
      \end{block}
    \end{column}
  \end{columns}
  \vspace{10pt}
  \begin{itemize}
    \item $S$ = sustituciones, $D$ = eliminaciones, $I$ = inserciones, $N$ = total de unidades en referencia
    \item Calculadas con \lib{jiwer.cer()} y \lib{jiwer.wer()} contra el \textit{ground truth}
    \item Complementadas con IoU (detección), SSIM y FID (calidad de inpainting)
  \end{itemize}
\end{frame}

% ══════════════════════════════════════════════════════════════════════════════
\section{Metodología}
% ══════════════════════════════════════════════════════════════════════════════

\begin{frame}{Visión General del Pipeline}
  \centering
  \begin{tikzpicture}[
    node distance=0.5cm,
    box/.style={rectangle, rounded corners=4pt, minimum width=2.1cm,
                minimum height=0.9cm, align=center, font=\scriptsize\bfseries,
                text=white, inner sep=4pt},
    arr/.style={->, thick, draw=MidGray}
  ]
    \node[box, fill=AccBlue]    (f0) {FASE 0\\Dataset};
    \node[box, fill=AccGreen,   right=of f0] (f1) {FASE 1\\Preproc.};
    \node[box, fill=AccAmber,   right=of f1] (f2) {FASE 2\\Detección};
    \node[box, fill=PriRed,     right=of f2] (f3) {FASE 3\\Inpainting};
    \node[box, fill=AccPurple,  right=of f3] (f4) {FASE 4\\OCR};
    \node[box, fill=AccTeal,    right=of f4] (f5) {FASE 5\\Validación};
    \draw[arr] (f0) -- (f1);
    \draw[arr] (f1) -- (f2);
    \draw[arr] (f2) -- (f3);
    \draw[arr] (f3) -- (f4);
    \draw[arr] (f4) -- (f5);
    \draw[arr, dashed, bend right=40, draw=AccAmber]
      (f2.south) to node[below, font=\tiny, text=AccAmber]{Sin sello} (f4.south);
    \node[below=0.15cm of f0, font=\tiny, text=MidGray] {StaVer\\Stampa};
    \node[below=0.15cm of f1, font=\tiny, text=MidGray] {OpenCV\\Hough};
    \node[below=0.15cm of f2, font=\tiny, text=MidGray] {HSV\\YOLOv8};
    \node[below=0.15cm of f3, font=\tiny, text=MidGray] {LaMa\\PyTorch};
    \node[below=0.15cm of f4, font=\tiny, text=MidGray] {Paddle\\Tesseract};
    \node[below=0.15cm of f5, font=\tiny, text=MidGray] {CER/WER\\IoU};
  \end{tikzpicture}
  \vspace{14pt}
  \begin{columns}[T]
    \begin{column}{0.32\textwidth}
      \small\rojo{\textbf{Entrada:}}\\Documento escaneado (JPG, PNG, PDF) con sello oficial superpuesto.
    \end{column}
    \begin{column}{0.32\textwidth}
      \small\rojo{\textbf{Innovación:}}\\Separación cromática selectiva — solo se borra el sello, no el texto debajo.
    \end{column}
    \begin{column}{0.32\textwidth}
      \small\rojo{\textbf{Salida:}}\\Texto extraído en JSON + reporte de métricas CER/WER.
    \end{column}
  \end{columns}
\end{frame}

% ─── ORIGEN Y DIVISION DE DATOS ──────────────────────────────────────────────
\begin{frame}{Datos — Origen y División del Dataset}
  % ── Tabla de fuentes ──────────────────────────────────────────
  \scriptsize
  \begin{tabular}{@{}llll@{}}
    \toprule
    \textbf{Fuente} & \textbf{Tipo} & \textbf{Uso} & \textbf{Cantidad} \\
    \midrule
    MultiFinBen-SpanishOCR & Docs.\ base (HuggingFace)  & Grupo A — sintético  & $\sim$500 docs \\
    StaVer (Kaggle)        & Sellos + máscaras GT        & Grupos A y B         & 400 imágenes   \\
    Stampa (Mendeley)      & Sellos sueltos              & Grupo A — sintético  & 1\,557 sellos  \\
    StaVer real            & Docs.\ auténticamente sellados & Grupo B — real    & $\sim$200 docs \\
    Transcripción manual   & Ground truth humano         & Grupo B — evaluación & 25 docs        \\
    \bottomrule
  \end{tabular}

  \vspace{10pt}
  % ── Barra de división ─────────────────────────────────────────
  \begin{tikzpicture}[font=\scriptsize]
    \fill[AccBlue!85]  (0,0)    rectangle (7.0, 0.6);
    \fill[AccAmber!85] (7.0,0)  rectangle (8.75, 0.6);
    \fill[PriRed!85]   (8.75,0) rectangle (10.5, 0.6);
    \node[white, font=\bfseries\scriptsize] at (3.5,  0.3) {Entrenamiento — 70\%};
    \node[white, font=\bfseries\scriptsize] at (7.875,0.3) {Validación — 15\%};
    \node[white, font=\bfseries\scriptsize] at (9.625,0.3) {Test — 15\%};
  \end{tikzpicture}

  \vspace{8pt}
  % ── Descripción grupos ────────────────────────────────────────
  \begin{columns}[T]
    \begin{column}{0.48\textwidth}
      \begin{block}{Grupo A — Sintético}
        \scriptsize Composición alfa (escala, rotación, opacidad aleatorias).
        Genera automáticamente imagen sucia, máscara y texto GT.
      \end{block}
    \end{column}
    \begin{column}{0.48\textwidth}
      \begin{block}{Grupo B — Real}
        \scriptsize Documentos auténticos de StaVer con sellos reales.
        25 transcritos manualmente — evaluación más rigurosa (Fase 5B).
      \end{block}
    \end{column}
  \end{columns}
\end{frame}

% ─── FASE 0 ───────────────────────────────────────────────────────────────────
\begin{frame}{Fase 0 — Construcción del Dataset}
  \begin{columns}[T]
    \begin{column}{0.48\textwidth}
      \begin{block}{Grupo A — Sintético}
        \begin{itemize}
          \item Documentos base: \textbf{MultiFinBen-SpanishOCR} (Hugging Face)
          \item Sellos: \textbf{StaVer} (Kaggle, 400 imgs.) + \textbf{Stampa} (Mendeley, 1\,557 sellos)
          \item Composición alfa con rotación, escala y opacidad aleatorias
          \item Salida automática: imagen sucia + máscara + ground truth textual
        \end{itemize}
      \end{block}
    \end{column}
    \begin{column}{0.48\textwidth}
      \begin{block}{Grupo B — Real}
        \begin{itemize}
          \item Documentos genuinamente sellados y escaneados del dataset StaVer
          \item 25 documentos transcritos \textbf{manualmente} para ground truth humano
          \item Usados en la Fase 5B — evaluación más rigurosa del proyecto
        \end{itemize}
      \end{block}
    \end{column}
  \end{columns}
\end{frame}

% ─── FASE 1 ───────────────────────────────────────────────────────────────────
\begin{frame}{Fase 1 — Preprocesamiento (1/2)}
  \textbf{Entrada:} imagen RGB escaneada \quad $\rightarrow$ \quad
  \textbf{Salida:} imagen binarizada + imagen a color corregida
  \vspace{10pt}
  \begin{enumerate}[label=\textbf{\rojo{\arabic*.}}, leftmargin=1.8em, itemsep=10pt]
    \item \textbf{Corrección de rotación}\\
      Transformada de Hough Probabilística detecta líneas dominantes. Cálculo del ángulo de inclinación promedio y rotación afín con interpolación bilineal.\\
      \lib{cv2.HoughLinesP} · \lib{cv2.getRotationMatrix2D} · \lib{cv2.warpAffine}
    \item \textbf{Eliminación de ruido}\\
      Filtro Non-Local Means ($h{=}10$, ventana $21{\times}21$) seguido de filtro Bilateral ($d{=}9$) que suaviza preservando los bordes del texto.\\
      \lib{cv2.fastNlMeansDenoising} · \lib{cv2.bilateralFilter}
  \end{enumerate}
\end{frame}

\begin{frame}{Fase 1 — Preprocesamiento (2/2)}
  \begin{enumerate}[label=\textbf{\rojo{\arabic*.}}, leftmargin=1.8em, itemsep=10pt]
    \setcounter{enumi}{2}
    \item \textbf{Binarización Adaptativa Gaussiana}\\
      Bloques $11{\times}11$, constante $C{=}2$. Cierre morfológico (kernel $2{\times}2$) para reconectar trazos fragmentados.\\
      \lib{cv2.adaptiveThreshold} · \lib{cv2.morphologyEx(MORPH\_CLOSE)}
  \end{enumerate}
\end{frame}

% ─── FASE 2 ───────────────────────────────────────────────────────────────────
\begin{frame}{Fase 2 — Detección del Sello}
  \textbf{Entrada:} imagen a color corregida \quad $\rightarrow$ \quad
  \textbf{Salida:} bounding box $(x, y, w, h)$ + máscara del sello
  \vspace{8pt}
  \begin{columns}[T]
    \begin{column}{0.48\textwidth}
      \begin{redblock}{\rojo{Ruta A — Segmentación HSV (baseline clásico)}}
        \begin{enumerate}[label=\arabic*., itemsep=3pt, leftmargin=1.2em]
          \item Conversión BGR $\rightarrow$ HSV
          \item Umbrales: Rojo $H \in [0°,10°]\cup[170°,180°]$,\\ Azul $H \in [100°,130°]$, $S > 0.4$
          \item Morfología: erosión $3{\times}3$, dilatación $5{\times}5$, cierre
          \item Filtrado de contornos por área y circularidad
        \end{enumerate}
        \lib{cv2.inRange} · \lib{cv2.morphologyEx} · \lib{cv2.findContours}
      \end{redblock}
    \end{column}
    \begin{column}{0.48\textwidth}
      \begin{blueblock}{\azul{Ruta B — YOLOv8-nano (deep learning)}}
        \begin{enumerate}[label=\arabic*., itemsep=3pt, leftmargin=1.2em]
          \item Fine-tuning desde COCO con transfer learning
          \item Entrenamiento: Grupo A + anotaciones Grupo B
          \item Inferencia con NMS ($\text{IoU} > 0.5$, confianza $\geq 0.5$)
          \item Salida: bounding box + máscara por instancia
        \end{enumerate}
        \lib{ultralytics.YOLO('yolov8n.pt')}
      \end{blueblock}
    \end{column}
  \end{columns}
  \vspace{8pt}
  \begin{alertblock}{Fase 2.5 — Clasificación del Sello}
    Saturación media $\bar{S}$ de la ROI en HSV.
    Si $\bar{S} > 0.25 \Rightarrow$ \textbf{COLOR} (Ruta 3A) \quad
    Si $\bar{S} \leq 0.25 \Rightarrow$ \textbf{NEGRO} (Ruta 3B)
  \end{alertblock}
\end{frame}

% ─── FASE 3 ───────────────────────────────────────────────────────────────────
\begin{frame}{Fase 3 — Eliminación del Sello con LaMa}
  \textbf{Entrada:} imagen a color + máscara + clasificación \quad $\rightarrow$ \quad
  \textbf{Salida:} \texttt{imagen\_limpia.png}
  \vspace{8pt}
  \begin{columns}[T]
    \begin{column}{0.48\textwidth}
      \begin{greenblock}{Ruta 3A — Sello de Color}
        \textbf{Separación cromática selectiva:}
        \begin{itemize}
          \item Cada píxel: \textbf{Texto} ($S < 0.2$, $V < 0.4$), \textbf{Sello} ($S > 0.35$) o \textbf{Mixto}
          \item LaMa reconstruye \emph{únicamente} píxeles de Sello y Mixto
          \item El texto visible se \textbf{preserva intacto}
        \end{itemize}
      \end{greenblock}
    \end{column}
    \begin{column}{0.48\textwidth}
      \begin{redblock}{Ruta 3B — Sello Negro}
        \textbf{Inpainting completo:}
        \begin{itemize}
          \item LaMa procesa toda la región de la máscara
          \item Reconstruye fondo y patrones periódicos (líneas)
          \item \alert{Limitación:} texto cubierto completamente no es recuperable
        \end{itemize}
      \end{redblock}
    \end{column}
  \end{columns}
  \vspace{6pt}
  \begin{exampleblock}{¿Por qué LaMa?}
    Convoluciones de Fourier (FFT) — campo receptivo \emph{global} desde la primera capa.
    Supera métodos locales en FID y LPIPS para máscaras de gran tamaño.\\
    \lib{simple-lama-inpainting} · \lib{torch} · GPU requerida
  \end{exampleblock}
\end{frame}

% ─── FASE 4 ───────────────────────────────────────────────────────────────────
\begin{frame}{Fase 4 — OCR y Post-corrección (1/2)}
  \textbf{Entrada:} imagen limpia \quad $\rightarrow$ \quad
  \textbf{Salida:} JSON con texto, coordenadas y scores de confianza
  \vspace{10pt}
  \begin{enumerate}[label=\textbf{\rojo{\arabic*.}}, leftmargin=1.8em, itemsep=10pt]
    \item \textbf{Motor principal — PaddleOCR v4}\\
      DBNet detecta polígonos de texto; SVTR (Vision Transformer) transcribe cada región. Configurado para español.\\
      \lib{PaddleOCR(lang='es', use\_angle\_cls=True)}
    \item \textbf{Decisión de confianza}\\
      Si confianza de un bloque $< 0.6$ → \textbf{Tesseract 5} (LSTM bidireccional) como motor de respaldo automático para ese bloque.\\
      \lib{pytesseract.image\_to\_string(lang='spa', config='--psm 6')}
  \end{enumerate}
\end{frame}

\begin{frame}{Fase 4 — OCR y Post-corrección (2/2)}
  \begin{enumerate}[label=\textbf{\rojo{\arabic*.}}, leftmargin=1.8em, itemsep=10pt]
    \setcounter{enumi}{2}
    \item \textbf{Post-corrección}\\
      Normalización de fechas y montos con \lib{re.sub}. Corrección ortográfica con \lib{pyspellchecker} (Levenshtein $\leq 2$). Detección de entidades (NER) con \lib{spaCy es\_core\_news\_sm} para validar nombres, fechas y montos.
  \end{enumerate}
\end{frame}

% ─── FASE 5 ───────────────────────────────────────────────────────────────────
\begin{frame}{Fase 5 — Validación y Métricas}
  \begin{table}
    \centering
    \small
    \begin{tabular}{@{}clll@{}}
      \toprule
      \textbf{Exp.} & \textbf{Tipo} & \textbf{Métrica} & \textbf{Descripción} \\
      \midrule
      \badge{AccBlue}{5A} & Automática   & CER · WER  & Pipeline vs.\ ground truth del Grupo A (sintético)           \\[4pt]
      \badge{PriRed}{5B}  & Manual       & CER · WER  & 25 docs.\ reales vs.\ transcripción humana                   \\[4pt]
      \badge{AccAmber}{5C}& Detección    & IoU medio  & HSV vs.\ YOLOv8 contra máscaras ground truth                 \\[4pt]
      \badge{AccTeal}{5D} & Inpainting   & SSIM · FID & Imagen restaurada vs.\ original limpio del Grupo A           \\
      \bottomrule
    \end{tabular}
  \end{table}
  \vspace{10pt}
  \begin{alertblock}{Hipótesis central}
    La eliminación de sellos reduce el CER en \textbf{al menos 5 puntos porcentuales} respecto al baseline sin procesamiento.
  \end{alertblock}
  \begin{footnotesize}
    Tabla comparativa de 4 condiciones: baseline crudo, pipeline HSV, pipeline YOLO, pipeline completo con post-corrección.
  \end{footnotesize}
\end{frame}

% ══════════════════════════════════════════════════════════════════════════════
\section{Innovación y Planificación}
% ══════════════════════════════════════════════════════════════════════════════

\begin{frame}{Mapeo de Tecnologías al Pipeline}
  \begin{table}
    \centering
    \scriptsize
    \begin{tabular}{@{}p{3.5cm}p{2.5cm}p{6.5cm}@{}}
      \toprule
      \textbf{Tecnología} & \textbf{Fase} & \textbf{Rol específico} \\
      \midrule
      OpenCV — Hough + Otsu           & Fase 1             & Corrección de rotación y binarización adaptativa \\
      HSV Thresholding                & Fase 2A            & Baseline clásico de detección de sellos por color \\
      YOLOv8-nano                     & Fase 2B            & Detección aprendida — comparada contra HSV por IoU \\
      Análisis HSV (saturación media) & Fase 2.5           & Clasifica el sello como COLOR o NEGRO \\
      LaMa — Separación cromática     & Fase 3A            & Elimina solo píxeles del sello, preserva texto \\
      LaMa — Inpainting completo      & Fase 3B            & Reconstruye la región completa del sello negro \\
      PaddleOCR v4                    & Fase 4 (principal) & DBNet + SVTR, configurado para español \\
      Tesseract 5                     & Fase 4 (fallback)  & Activado cuando confianza del bloque $< 0.6$ \\
      pyspellchecker + spaCy          & Fase 4.5           & Corrección ortográfica y NER para validar entidades \\
      jiwer (CER/WER) + SSIM/FID      & Fase 5             & Medición cuantitativa del impacto del pipeline \\
      \bottomrule
    \end{tabular}
  \end{table}
\end{frame}

% ─── CIERRE ──────────────────────────────────────────────────────────────────
\begin{frame}[standout]
  \centering
  {\Large\bfseries Gracias}\\[16pt]
  {\normalsize Detección y Eliminación Inteligente de Sellos\\para Mejora del OCR en Documentos Oficiales}\\[20pt]
  {\small Manuel Cruz Garrote}\\
  {\small Universidad del Rosario · Marzo 2026}
\end{frame}

% ─── REFERENCIAS ─────────────────────────────────────────────────────────────
\begin{frame}[allowframebreaks]{Referencias}
  \scriptsize
  \begin{enumerate}[label={[\arabic*]}, leftmargin=1.8em, itemsep=4pt]
    \item Suvorov, R. et al. (2022). Resolution-robust Large Mask Inpainting with Fourier Convolutions. \textit{WACV 2022}.
    \item Du, Y. et al. (2023). PP-OCRv4: Self-Supervised Pre-Training for Text Detection and Recognition. \textit{arXiv:2309.06383}.
    \item Jocher, G. et al. (2023). Ultralytics YOLOv8. \textit{GitHub}.
    \item Ravi, N. et al. (2024). SAM 2: Segment Anything in Images and Videos. \textit{Meta AI, arXiv:2408.00714}.
    \item Liang, J. et al. (2022). An adaptive deskewing algorithm for document images. \textit{Pattern Recognition Letters}.
    \item Terven, J. \& Cordova, D. (2023). A Comprehensive Review of YOLO Architectures. \textit{arXiv:2304.00501}.
    \item Smith, R. (2021+). Tesseract OCR Engine v5. \textit{GitHub}.
  \end{enumerate}
\end{frame}

\end{document}
